\section{PbI\texorpdfstring{$_2$}{2} stacking-disorder model}
\label{sec:si_pbi2_transition}

This section gives the PbI$_2$-specific state construction and the finite-stack
intensity used in the main text.  The general strategy--propagating layer-pair
probabilities with a transition matrix--follows the established theory of
one-dimensionally disordered crystals.\cite{HendricksTeller1942,KakinokiKomura1954,KakinokiKomura1965,Hart2018}
The particular six states, allowed interface events, and fault templates below
are the model assumptions introduced for the present PbI$_2$ analysis.

\subsection{Trilayer amplitudes and six-state basis}

As in the main text, $\rho\in\{0,1,2\}$ labels the basal registry and $\sigma\in\{+,-\}$ labels trilayer orientation. Here $\sigma$ is a discrete state label; the mosaic width $\sigma$ in Sec.~\ref{sec:si_mosaic} is a separate, continuous parameter. The basal registries are
\begin{equation}
  0=(0,0),\qquad
  1=\left(\frac13,\frac23\right),\qquad
  2=\left(\frac23,\frac13\right).
  \label{eq:si_pbi2_registries}
\end{equation}
Taking the Pb plane as the local origin and writing the Pb--I vertical offset
as $\eta=\Delta z_{\mathrm{I-Pb}}/c$, the two orientation-dependent trilayer
amplitudes are
\begin{align}
 F_+(h,k,L)={}&f_{\mathrm{Pb}}+
 f_{\mathrm I}\exp\!\left\{2\pi i\left[\frac{2h+k}{3}+L\eta\right]\right\}
 +f_{\mathrm I}\exp\!\left\{2\pi i\left[\frac{h+2k}{3}-L\eta\right]\right\},
 \label{eq:si_pbi2_Fplus}\\
 F_-(h,k,L)={}&f_{\mathrm{Pb}}+
 f_{\mathrm I}\exp\!\left\{2\pi i\left[\frac{2h+k}{3}-L\eta\right]\right\}
 +f_{\mathrm I}\exp\!\left\{2\pi i\left[\frac{h+2k}{3}+L\eta\right]\right\}.
 \label{eq:si_pbi2_Fminus}
\end{align}
Atomic form factors retain their dependence on $Q=|\mathbf G_{hk}(L)|$ and source wavelength $\lambda_s$. These arguments are suppressed in the matrix expressions below.  The two amplitudes
satisfy $F_-(h,k,L)=F_+(h,k,-L)$.

For the forward registry step $\mathbf{s}=(1/3,2/3)$, define
\begin{equation}
  \omega_{hk}=\exp\!\left[2\pi i(hs_x+ks_y)\right]
  =\exp\!\left[\frac{2\pi i}{3}(h+2k)\right].
  \label{eq:si_pbi2_omega}
\end{equation}
The ordered state basis and its scattering-amplitude vector are
\begin{equation}
 \mathcal{S}=[0F_+,1F_+,2F_+,0F_-,1F_-,2F_-],\qquad
 \mathbf{g}=\begin{pmatrix}
 F_+&\omega_{hk} F_+&\omega_{hk}^2F_+&F_-&\omega_{hk} F_-&\omega_{hk}^2F_-
 \end{pmatrix}^{\mathsf T}.
 \label{eq:si_pbi2_state_amplitudes}
\end{equation}
Because $\omega_{hk}^3=1$, registry 2 is also one backward step from registry 0.

\subsection{Direction-resolved interface law}

The five mutually exclusive interface events are
\begin{equation}
 \boldsymbol{\theta}=(a,b_+,b_-,d_+,d_-),\qquad
 a+b_++b_-+d_++d_-=1.
 \label{eq:si_pbi2_interface_law}
\end{equation}
The entries of $\boldsymbol{\theta}$ are event probabilities. Here $a$ is the probability of retaining both registry and orientation; the basal lattice constant denoted by $a$ in Sec.~\ref{sec:si_structure_factor_conventions} is a different quantity. The $a$ event preserves registry and orientation. The $b_+$ and $b_-$ events shift the
registry by $+1$ and $-1$ while preserving orientation. The $d_+$ and $d_-$ events
combine an orientation change with the two handed shift rules.  Every allowed orientation flip includes a registry shift.  For a $d_+$ event, an $F_+\rightarrow F_-$ step
advances the registry and the corresponding $F_-\rightarrow F_+$ step returns
it.  The directions are interchanged for $d_-$.  This reversal is what makes a
pure $d_+$ or $d_-$ process a deterministic two-layer 4H repeat.

Let
\begin{equation}
 S_+=\begin{pmatrix}0&1&0\\0&0&1\\1&0&0\end{pmatrix},\qquad
 S_-=S_+^{-1}=S_+^{\mathsf T}.
 \label{eq:si_pbi2_shift_matrices}
\end{equation}
The matrices $S_+$ and $S_-$ shift the registry forward and backward, respectively. Let $I_m$ denote the $m$-dimensional identity matrix. With rows denoting the current state and columns the next state, define
\begin{align}
 A&=aI_3+b_+S_+ +b_-S_-,\\
 B&=d_+S_+ +d_-S_-,\\
 C&=d_+S_- +d_-S_+.
\end{align}
Here $A$ is the orientation-preserving block, and $B$ and $C$ are the two orientation-changing blocks. The full six-state transition matrix is
\begin{equation}
 T_6=\begin{pmatrix}A&B\\C&A\end{pmatrix}.
 \label{eq:si_pbi2_T6}
\end{equation}
For example, the row for the state $0F_+$ is
$(a,b_+,b_-,0,d_+,d_-)$.  Every row of $T_6$ sums to one.

The deterministic parent-event vectors, in the order
$(a,b_+,b_-,d_+,d_-)$, are
\begin{equation}
 \mathbf e_{2H}=(1,0,0,0,0),\qquad
 \mathbf e_{4H}=(0,0,0,1,0),\qquad
 \mathbf e_{6H}=(0,1,0,0,0).
 \label{eq:si_pbi2_parent_vectors}
\end{equation}
The opposite-handed 4H and 6H parents are obtained by selecting $d_-$ and
$b_-$ instead.

\subsection{Exact rod-dependent reduction}

The three registry-shift matrices are diagonal in the discrete Fourier basis.
In the sector selected by the scattering phase $\omega_{hk}$, the $6\times6$
transition reduces exactly to
\begin{equation}
 M(\omega_{hk})=
 \begin{pmatrix}
 a+b_+\omega_{hk}+b_-\omega_{hk}^{-1} & d_+\omega_{hk}+d_-\omega_{hk}^{-1}\\
 d_+\omega_{hk}^{-1}+d_-\omega_{hk} & a+b_+\omega_{hk}+b_-\omega_{hk}^{-1}
 \end{pmatrix}.
 \label{eq:si_pbi2_Momega}
\end{equation}
All registry states are retained. Equation~\eqref{eq:si_pbi2_Momega} is one Fourier
block of Eq.~\eqref{eq:si_pbi2_T6}.  Setting the registry phase to unity gives
the orientation-population matrix
\begin{equation}
 P=M(1)=
 \begin{pmatrix}
 a+b_++b_- & d_++d_-\\
 d_++d_- & a+b_++b_-
 \end{pmatrix}.
 \label{eq:si_pbi2_P}
\end{equation}
The matrix $M(\omega_{hk})$ carries complex scattering phases, while $P$ is a row-stochastic probability matrix.

\subsection{Finite-stack intensity}

Let $d$ be the separation between neighboring trilayer centers and
$\xi=\exp(iq_zd)$ the corresponding vertical phase, with $q_z=2\pi L/c$.  Define
\begin{equation}
 \mathbf f=\begin{pmatrix}F_+\\F_-\end{pmatrix},\qquad
 D_{F^*}=\operatorname{diag}(F_+^*,F_-^*).
 \label{eq:si_pbi2_form_vector}
\end{equation}
If $\boldsymbol{\pi}^{\mathrm{ori}}_0$ is the initial orientation row vector,
then the orientation population at trilayer $n$ is
\begin{equation}
 \boldsymbol{\pi}^{\mathrm{ori}}_n
 =\boldsymbol{\pi}^{\mathrm{ori}}_0P^n.
 \label{eq:si_pbi2_orientation_population}
\end{equation}
The vector $\mathbf f$ collects the two orientation amplitudes, and $D_{F^*}$ is their diagonal complex-conjugate matrix. The superscript $\mathrm{ori}$ labels orientation probabilities; each row vector $\boldsymbol{\pi}^{\mathrm{ori}}_n$ sums to one. For integer layer separation $1\leq\Delta<N$, the phase-weighted correlation between trilayers $n$ and $n+\Delta$ is
\begin{equation}
 K_{n,\Delta}=\boldsymbol{\pi}^{\mathrm{ori}}_n
 D_{F^*}M(\omega_{hk})^\Delta\mathbf f.
 \label{eq:si_pbi2_pair_kernel}
\end{equation}
For one fixed parent-rich population $j\in\{2H,4H,6H\}$, the main-text rod intensity uses the per-trilayer convention $S_{hk}^{(j)}=\langle I_N^{(j)}\rangle/N$, where $I_N^{(j)}$ is the intensity of the complete $N$-trilayer stack before the disorder average. Population labels on $F_\pm$, the matrices, and $K_{n,\Delta}$ are suppressed in this single-population derivation. For a finite stack of $N$ trilayers,
\begin{align}
 S_{hk}^{(j)}(L,\lambda_s)={}&
 \frac{1}{N}\sum_{n=0}^{N-1}
 \boldsymbol{\pi}^{\mathrm{ori}}_n
 \begin{pmatrix}|F_+|^2\\|F_-|^2\end{pmatrix}
 \nonumber\\
 &+\frac{2}{N}\operatorname{Re}\!\left[
 \sum_{\Delta=1}^{N-1}\xi^\Delta
 \sum_{n=0}^{N-1-\Delta}K_{n,\Delta}\right].
 \label{eq:si_pbi2_finite_intensity}
\end{align}
The connection to the main-text pair-correlation notation is
\begin{equation}
 C_{hk}^{(j)}(\Delta;L,\lambda_s)
 =\frac{1}{N-\Delta}\sum_{n=0}^{N-1-\Delta}K_{n,\Delta},
 \qquad 1\leq\Delta<N.
 \label{eq:si_pbi2_main_correlation}
\end{equation}
The first term in Eq.~\eqref{eq:si_pbi2_finite_intensity} is $S_{\mathrm{self},hk}^{(j)}(L,\lambda_s)$. Since $\xi^\Delta=e^{iq_zd\Delta}$, grouping the pair sum with this definition gives the main-text finite-stack expression directly.

The first term contains the $N$ self contributions and the second contains all
distinct trilayer pairs.  The calculation evaluates these finite sums
directly, retaining end effects and avoiding singular closed forms when an
eigenvalue approaches unity.  Equivalent matrix-power and recursive
formulations are standard for planar-fault diffraction.\cite{Treacy1991}

\subsection{Polytype-rich fault templates and population mixture}

For each parent-rich population,
\begin{equation}
 \boldsymbol{\theta}_j(\epsilon_j)
 =(1-\epsilon_j)\mathbf e_j+\epsilon_j\mathbf q_j,
 \qquad j\in\{2H,4H,6H\}.
 \label{eq:si_pbi2_fault_parameterization}
\end{equation}
Here $\mathbf e_j$ is the deterministic parent-event vector, $\mathbf q_j$ is the probability vector for a fault event, and $0\leq\epsilon_j\leq1$ is its probability. The equal-alternative templates used in the present implementation are
\begin{align}
 \mathbf q_{2H}&=\left(0,\frac14,\frac14,\frac14,\frac14\right),\\
 \mathbf q_{4H}&=\left(\frac14,\frac14,\frac14,0,\frac14\right),\\
 \mathbf q_{6H}&=\left(\frac14,0,\frac14,\frac14,\frac14\right).
 \label{eq:si_pbi2_fault_templates}
\end{align}
Consequently, $\epsilon_j$ is exactly the probability that an interface in
population $j$ follows one of the four alternative events. Its numerical value is conditional on the chosen $\mathbf q_j$ and is interpreted together with that template.

The parent-rich populations are mutually incoherent, with nonnegative weights $w_j$ that sum to one. Their structural intensities combine using the main-text notation:
\begin{equation}
 S_{hk}(L,\lambda_s)=\sum_{j\in\{2H,4H,6H\}}
 w_jS_{hk}^{(j)}(L,\lambda_s),\qquad
 w_{2H}+w_{4H}+w_{6H}=1.
 \label{eq:si_pbi2_total_intensity}
\end{equation}
The detector/profile prediction, background, scale and optional finite-profile kernel are specified separately in Sec.~\ref{subsec:si_detector_derived_objectives}.
The core parameterization has three disorder magnitudes and two independent
population weights.  The stack length $N$, scale, and any fixed or refined
finite-profile and instrumental broadening terms are separate.  When three independent disorder magnitudes remain underconstrained, the constrained starting model is
$\epsilon_{2H}=\epsilon_{4H}=\epsilon_{6H}$.

\subsection{Deterministic and phase-insensitive checks}

Table~\ref{tab:si_parent_cycles} gives the zero-disorder cycles.
\begin{table}[!htbp]
\centering\small
\caption{Deterministic PbI$_2$ parent cycles from the initial state $0F_+$. Each row specifies the unit-probability interface event, its registry/orientation cycle, and the corresponding reduced transfer matrix. The phase $\omega_{hk}$ is defined in Eq.~\eqref{eq:si_pbi2_omega}; $I_2$ is the two-dimensional identity matrix.}
\label{tab:si_parent_cycles}
\begin{tabular}{llll}
\toprule
Parent & Nonzero event & State cycle from $0F_+$ & Reduced matrix\\
\midrule
2H & $a=1$ & $0F_+\rightarrow0F_+$ & $I_2$\\
4H & $d_+=1$ & $0F_+\rightarrow1F_-\rightarrow0F_+$ &
$\left(\begin{smallmatrix}0&\omega_{hk}\\\omega_{hk}^{-1}&0\end{smallmatrix}\right)$\\
6H & $b_+=1$ & $0F_+\rightarrow1F_+\rightarrow2F_+\rightarrow0F_+$ &
$\omega_{hk} I_2$\\
\bottomrule
\end{tabular}
\end{table}
\FloatBarrier
For 4H, $M(\omega_{hk})^2=I_2$. For 6H, $M(\omega_{hk})^3=I_2$ because
$\omega_{hk}^3=1$.  The $d_-=1$ and $b_-=1$ limits generate the opposite-handed
cycles.

For $h=k=0$, $\omega_{hk}=1$ and $F_+=F_-=F$.  Because $P^\Delta\mathbf 1=\mathbf 1$, where $\mathbf 1=(1,1)^{\mathsf T}$, Eq.~\eqref{eq:si_pbi2_pair_kernel} becomes
$K_{n,\Delta}=|F|^2$ for every allowed transition law.  Equation
\eqref{eq:si_pbi2_finite_intensity} then reduces to the normalized Laue factor
\begin{equation}
 S_{00}^{(j)}(L,\lambda_s)
 =\frac{|F|^2}{N}\left|\sum_{n=0}^{N-1}\xi^n\right|^2
 =\frac{|F|^2}{N}\frac{\sin^2(N\vartheta/2)}{\sin^2(\vartheta/2)},
 \qquad \vartheta=q_zd,\quad \xi=e^{i\vartheta}.
 \label{eq:si_pbi2_r0_laue}
\end{equation}
Thus the strong $r=0$ oscillations are finite-thickness fringes. The included registry and orientation paths are contrast-matched on this rod, so faults can remain present with the same axial intensity. Figure~\ref{fig:si_pbi2_transition_validation} reserves the direct numerical checks of these limits. Figure~\ref{fig:si_pbi2_sensitivity} separates the effects of faults within a stack from changes in population weights.


\begin{figure}[!htbp]
  \centering
  \siplaceholder{PbI$_2$ model verification}{(a) Full six-state and reduced calculations.\\(b) Deterministic 2H, 4H, and 6H limits.\\(c) Axial invariance across transition laws.}
  \caption{Placeholder for verification of the PbI$_2$ model. Direct arrays will compare the full and reduced expressions, independent deterministic-parent sums, and the phase-insensitive axial limit.}
  \label{fig:si_pbi2_transition_validation}
\end{figure}


\begin{figure}[!htbp]
  \centering
  \siplaceholder{PbI$_2$ disorder and population sensitivity}{(a) One fault probability varied at fixed population weights.\\(b) Population weights varied at fixed transition laws.\\Shared axes and declared structural assumptions.}
  \caption{Placeholder for direct PbI$_2$ sensitivity calculations. Separate controlled panels distinguish within-stack fault probabilities from incoherent population weights.}
  \label{fig:si_pbi2_sensitivity}
\end{figure}
